\title[Synthesis and exam preparation]{Synthesis, design showcase and exam preparation}
\section{Synthesis and exam preparation}

\begin{frame}[plain]
	\titlepage
\end{frame}

%%%%%%%%%%%%%%%%%%%%%%%%%%%%%%%%%%%%%%%%%%%%%%%%%%%%%%%%%%%%%
%% Overview %%
%%%%%%%%%%%%%%%%%%%%%%%%%%%%%%%%%%%%%%%%%%%%%%%%%%%%%%%%%%%%%
\subsection{Synthesis}

\begin{frame}
	\frametitle{Overview}
	Fourteen lectures took the converter apart: devices, layout, drivers, magnetics, topologies, control, EMI and heat. Today we put it back together and prepare the exam.
	\begin{itemize}
		\item The frequency ladder revisited: what limits each rung, and what moves the limit
		\item One design, five questions
		\item The oral exam: format, assessment, example questions
		\item Design showcase: three-minute pitches with peer questions
		\item Course evaluation, theses and projects
	\end{itemize}
	After this session you know how the exam works and can present your design in three minutes, with the number that matters most for each decision.
\end{frame}

%%%%%%%%%%%%%%%%%%%%%%%%%%%%%%%%%%%%%%%%%%%%%%%%%%%%%%%%%%%%%
%% Ladder revisited %%
%%%%%%%%%%%%%%%%%%%%%%%%%%%%%%%%%%%%%%%%%%%%%%%%%%%%%%%%%%%%%
\begin{frame}
	\frametitle{The ladder revisited}
	\begin{figure}
		\centering
		\begin{tikzpicture}[
				rung/.style={draw, rounded corners=2pt, minimum width=1.7cm, minimum height=0.75cm, font=\small\bfseries},
				limit/.style={draw, rounded corners=2pt, text width=3.6cm, align=left, font=\scriptsize, minimum height=0.8cm, inner sep=3pt},
				answer/.style={draw, rounded corners=2pt, text width=5.9cm, align=left, font=\scriptsize, minimum height=0.8cm, inner sep=3pt},
			]
			\draw[line width=1.5pt, gray] (-1.1,-0.5) -- (-1.1,4.4);
			\draw[line width=1.5pt, gray] (1.1,-0.5) -- (1.1,4.4);
			\foreach \y/\f in {0/{100\,kHz}, 1.3/{1\,MHz}, 2.6/{10\,MHz}, 3.9/{100\,MHz}} {
					\node[rung, fill=white] at (0,\y) {\f};
				}
			\node[limit, anchor=west] at (1.6,0) {\textbf{Switching loss}};
			\node[limit, anchor=west] at (1.6,1.3) {\textbf{Parasitics and timing}};
			\node[limit, anchor=west] at (1.6,2.6) {\textbf{Magnetics and $C_\mathrm{oss}$ loss}};
			\node[limit, anchor=west] at (1.6,3.9) {\textbf{Packaging, control, heat}};
			\node[answer, anchor=west] at (5.6,0) {WBG devices, $Q_\mathrm{oss}V$ instead of $E_\mathrm{oss}$, ZVS by charge in the dead time, TCM, resonant tanks (L2, L8, L9)};
			\node[answer, anchor=west] at (5.6,1.3) {vertical loops, Kelvin source, damped gate loops, ns timing budget, probes four times faster than the edge (L3--L5)};
			\node[answer, anchor=west] at (5.6,2.6) {performance factor, low-$\mu$ and air-core magnetics, single-switch inverters, SC and hybrid converters (L6, L7, L10, L12)};
			\node[answer, anchor=west] at (5.6,3.9) {on/off and ripple-based control, integration, designed filters and thermal paths (L13, L14)};
			\node[font=\scriptsize\bfseries, anchor=west] at (1.6,4.6) {What breaks};
			\node[font=\scriptsize\bfseries, anchor=west] at (5.6,4.6) {What the course offers};
			\draw[-{Latex[length=3mm]}, thick] (-1.8,-0.4) -- node[sloped, above, font=\scriptsize] {$f_\mathrm{sw}\times 10$ per rung} (-1.8,4.4);
		\end{tikzpicture}
		\caption{The frequency ladder of L1 with the answers developed in L2 to L14}
	\end{figure}
\end{frame}

%%%%%%%%%%%%%%%%%%%%%%%%%%%%%%%%%%%%%%%%%%%%%%%%%%%%%%%%%%%%%
%% Limits today %%
%%%%%%%%%%%%%%%%%%%%%%%%%%%%%%%%%%%%%%%%%%%%%%%%%%%%%%%%%%%%%
\begin{frame}
	\frametitle{What limits the rungs today}
	\begin{columns}[T]
		\begin{column}{0.48\textwidth}
			\footnotesize
			\textbf{Devices}
			\begin{itemize}
				\item $C_\mathrm{oss}$ hysteresis loss in Si and GaN transistors and SiC diodes: soft switching is not loss-free (L2, L5, L10).
				\item Dynamic $R_\mathrm{DS(on)}$ of GaN after high-voltage blocking (L2).
			\end{itemize}
			\textbf{Magnetics}
			\begin{itemize}
				\item The performance factor of available materials improves up to about 30\,MHz; beyond, air core (L6, L10).
				\item Carrying high current at high frequency: the copper skin depth is 21\,$\mu$m at 10\,MHz, and litz wire loses its advantage above several MHz (L7).
			\end{itemize}
		\end{column}
		\begin{column}{0.48\textwidth}
			\footnotesize
			\textbf{Packaging and heat}
			\begin{itemize}
				\item Loop and common-source inductance move into the package; the heat density grows as the volume shrinks (L3, L14).
			\end{itemize}
			\textbf{Control and EMI}
			\begin{itemize}
				\item Delays of sensing and control are fixed in nanoseconds and do not shrink with the period (L13).
				\item Filters shrink with $f_\mathrm{sw}$ only if their parasitics are designed (L14).
			\end{itemize}
			\begin{block}{Key message}
				Each rung has a known limit and known tools; your design shows which limit you meet first.
			\end{block}
		\end{column}
	\end{columns}
	\blfootnote{Sources: Perreault, APEC 2024 plenary; A.\,J. Hanson et al., IEEE TPEL, vol.\ 31, no.\ 11, 2016.}
\end{frame}

%%%%%%%%%%%%%%%%%%%%%%%%%%%%%%%%%%%%%%%%%%%%%%%%%%%%%%%%%%%%%
%% Room at the top %%
%%%%%%%%%%%%%%%%%%%%%%%%%%%%%%%%%%%%%%%%%%%%%%%%%%%%%%%%%%%%%
\begin{frame}
	\frametitle{Plenty of room at the top}
	\begin{columns}[T]
		\begin{column}{0.48\textwidth}
			\small
			Perreault (APEC 2024): much higher frequencies appear possible, tens of MHz at kW power levels. The challenges are magnetics, devices, packaging, circuit design and control, the topics of this course.
			\vspace{0.2cm}

			Examples presented there:
			\begin{itemize}
				\item 5\,kW, 13.56\,MHz inverter with fast response and wide power range (H. Zhang, MIT, 2022)
				\item 570\,nH inductor at 13.56\,MHz, $Q=1150$ at 80\,A peak (M. Joisher et al., APEC 2024)
			\end{itemize}
		\end{column}
		\begin{column}{0.48\textwidth}
			\small
			\begin{itemize}
				\item Inductive power transfer: 1\,kW at 6.78\,MHz with 95\,\% (Stanford, 2021), 3\,kW at 3.39\,MHz with 95\,\% (Imperial College, 2023)
				\item Capacitive power transfer: 3.7\,kW at 13.56\,MHz with 93\,\% (Cornell, 2022)
				\item 30\,W LED driver at 5--10\,MHz and 250\,W PFC supply at 1--4\,MHz with small EMI filters (L14)
			\end{itemize}
			\begin{block}{Key message}
				High frequency is no longer a laboratory curiosity; it is a design choice with known costs.
			\end{block}
		\end{column}
	\end{columns}
	\blfootnote{Source: D.\,J. Perreault, plenary talk, IEEE APEC 2024, and the works cited there.}
\end{frame}

%%%%%%%%%%%%%%%%%%%%%%%%%%%%%%%%%%%%%%%%%%%%%%%%%%%%%%%%%%%%%
%% Five questions %%
%%%%%%%%%%%%%%%%%%%%%%%%%%%%%%%%%%%%%%%%%%%%%%%%%%%%%%%%%%%%%
\begin{frame}
	\frametitle{One design, five questions}
	\begin{table}
		\centering
		\footnotesize
		\begin{tabular}{|c|>{\raggedright\arraybackslash}p{4.6cm}|>{\raggedright\arraybackslash}p{6.4cm}|c|}
			\hline
			  & \textbf{Question}                                & \textbf{The number that answers it}                                                   & \textbf{Lectures} \\
			\hline
			1 & Where does the switching energy go?              & $Q_\mathrm{oss}V_\mathrm{DC}f_\mathrm{sw}$ hard, or the charge $2Q_\mathrm{oss}$ within the dead time & L2, L8      \\
			2 & Where is the energy stored, and what limits it?  & $\hat{B}$ from core loss, $F_\mathrm{R}$ of the winding, capacitance at DC bias        & L6, L7, L12       \\
			3 & Which parasitics are part of the circuit?        & $L_\mathrm{loop}$, $C_\mathrm{oss}$, winding and package capacitance in the tank        & L3, L9, L10       \\
			4 & How is it driven, controlled and measured?       & delays against $T_\mathrm{s}$, crossover and margin, probe bandwidth                  & L4, L5, L13       \\
			5 & Where do the heat and the noise go?              & $R_\mathrm{th}$ from junction to ambient, filter corner, CM capacitance              & L14               \\
			\hline
		\end{tabular}
		\caption{The five questions behind every design review and the oral exam}
	\end{table}
	\small
	$f_\mathrm{sw}$ is the answer to none of them alone: it is the variable that trades all five against volume.
\end{frame}

%%%%%%%%%%%%%%%%%%%%%%%%%%%%%%%%%%%%%%%%%%%%%%%%%%%%%%%%%%%%%
%% Oral exam %%
%%%%%%%%%%%%%%%%%%%%%%%%%%%%%%%%%%%%%%%%%%%%%%%%%%%%%%%%%%%%%
\subsection{Exam preparation}

\begin{frame}
	\frametitle{The oral exam}
	\begin{columns}[T]
		\begin{column}{0.56\textwidth}
			\begin{table}
				\centering
				\footnotesize
				\begin{tabular}{|l|c|>{\raggedright\arraybackslash}p{3.3cm}|}
					\hline
					\textbf{Part}   & \textbf{Time}    & \textbf{Content}                                     \\
					\hline
					A Presentation  & $\approx$10\,min & your design: at most 5 slides or one A3 design sheet \\
					B Interrogation & $\approx$20\,min & what-if questions with short calculations on paper   \\
					C Breadth       & $\approx$10\,min & 2--3 question chits on L1--L14                       \\
					\hline
				\end{tabular}
				\caption{Oral exam of about 40 minutes, individual, with an assessor}
			\end{table}
		\end{column}
		\begin{column}{0.41\textwidth}
			\small
			\begin{itemize}
				\item Each team member is examined alone and must be able to defend the whole design.
				\item Without an own design: the reference design announced in advance replaces it in parts A and B.
				\item Measured results are not required for a top grade; a sound validation plan is.
				\item Bring a calculator; the design package (M4) is the starting point.
			\end{itemize}
		\end{column}
	\end{columns}
\end{frame}

%%%%%%%%%%%%%%%%%%%%%%%%%%%%%%%%%%%%%%%%%%%%%%%%%%%%%%%%%%%%%
%% Assessment %%
%%%%%%%%%%%%%%%%%%%%%%%%%%%%%%%%%%%%%%%%%%%%%%%%%%%%%%%%%%%%%
\begin{frame}
	\frametitle{Assessment}
	\begin{table}
		\centering
		\footnotesize
		\begin{tabular}{|l|c|>{\raggedright\arraybackslash}p{7.6cm}|}
			\hline
			\textbf{Criterion}               & \textbf{Weight} & \textbf{What is assessed}                                                                     \\
			\hline
			Design reasoning and correctness & 30\,\%          & assumptions, calculations, a consistent loss budget, correct use of the models                \\
			Depth under what-if questions    & 25\,\%          & quantitative prediction of consequences, identification of the limiting mechanism             \\
			Breadth                          & 25\,\%          & chit questions across the course                                                              \\
			Engineering judgement            & 20\,\%          & trade-offs, validation and measurement plan, safety, limitations, critical self-assessment    \\
			\hline
		\end{tabular}
		\caption{Criteria of the oral exam; the grade is based on the oral performance only}
	\end{table}
	\small
	\textbf{How to prepare:} know the five numbers of your design by heart (efficiency, $f_\mathrm{sw}$, the dominant loss, the peak stress, the hot spot); recalculate one what-if per lecture; practise the chits in groups of three (Tutorial 14).
\end{frame}

%%%%%%%%%%%%%%%%%%%%%%%%%%%%%%%%%%%%%%%%%%%%%%%%%%%%%%%%%%%%%
%% Example questions %%
%%%%%%%%%%%%%%%%%%%%%%%%%%%%%%%%%%%%%%%%%%%%%%%%%%%%%%%%%%%%%
\begin{frame}
	\frametitle{Example questions}
	\begin{columns}[T]
		\begin{column}{0.48\textwidth}
			\small
			\textbf{Part B: what if \ldots}
			\begin{itemize}
				\item the load is halved: is ZVS kept?
				\item $f_\mathrm{sw}$ is doubled: what fails first?
				\item the loop inductance doubles: how large is the overshoot?
				\item the MLCC loses 60\,\% of its capacitance at bias?
				\item the core runs 20\,K hotter: loss and saturation margin?
				\item you have to trust an efficiency of 98\,\%: which measurement?
			\end{itemize}
		\end{column}
		\begin{column}{0.48\textwidth}
			\small
			\textbf{Part C: chits (examples)}
			\begin{itemize}
				\item Why is the hard turn-on loss $Q_\mathrm{oss}V_\mathrm{DC}$ and not $E_\mathrm{oss}$?
				\item Why does the optimum thickness of a winding layer fall with the number of layers?
				\item What limits the frequency of a class E inverter?
				\item Why does charging a capacitor from a voltage source lose $\frac12C\Delta V^2$, whatever the resistance?
				\item What does a delay of $1.5\,T_\mathrm{s}$ cost at $f_\mathrm{sw}/10$?
			\end{itemize}
		\end{column}
	\end{columns}
\end{frame}

%%%%%%%%%%%%%%%%%%%%%%%%%%%%%%%%%%%%%%%%%%%%%%%%%%%%%%%%%%%%%
%% Showcase format %%
%%%%%%%%%%%%%%%%%%%%%%%%%%%%%%%%%%%%%%%%%%%%%%%%%%%%%%%%%%%%%
\subsection{Design showcase}

\begin{frame}
	\frametitle{Design showcase}
	\begin{columns}[T]
		\begin{column}{0.48\textwidth}
			\small
			\textbf{Each team: 3 minutes pitch, 2 minutes questions}
			\begin{enumerate}
				\item Specification and topology, in one sentence
				\item $f_\mathrm{sw}$, efficiency and volume: the three numbers
				\item The decision you can defend best, with its number
				\item The largest risk and how you would test it
			\end{enumerate}
		\end{column}
		\begin{column}{0.48\textwidth}
			\small
			\textbf{Peer questions}
			\begin{itemize}
				\item Every team prepares one question for the team after it; the next slides give a starting point for each brief.
				\item Ask for a number, not an opinion.
			\end{itemize}
			\textbf{Demos} from the hardware track follow the pitches.
		\end{column}
	\end{columns}
\end{frame}

%%%%%%%%%%%%%%%%%%%%%%%%%%%%%%%%%%%%%%%%%%%%%%%%%%%%%%%%%%%%%
%% Showcase B1 %%
%%%%%%%%%%%%%%%%%%%%%%%%%%%%%%%%%%%%%%%%%%%%%%%%%%%%%%%%%%%%%
\begin{frame}
	\frametitle{B1: ZVS GaN buck}
	\begin{varblock}[\textwidth]{48\,V to 12\,V, 100\,W, 1--3\,MHz, triangular current mode (TCM) or quasi-square-wave (QSW)}
		\begin{itemize}
			\item Which negative current gives ZVS, and what does it cost in RMS current (L8)?
			\item What sets $f_\mathrm{sw}$ at light load, and how do you limit it?
			\item Where is the hot spot of the GaN transistor, and how many vias carry its heat (L14)?
		\end{itemize}
	\end{varblock}
	\small
	Numbers to show: $f_\mathrm{sw}$ at full and 10\,\% load, negative current, $I_\mathrm{rms}$, $T_\mathrm{j}$.
\end{frame}

%%%%%%%%%%%%%%%%%%%%%%%%%%%%%%%%%%%%%%%%%%%%%%%%%%%%%%%%%%%%%
%% Showcase B2 %%
%%%%%%%%%%%%%%%%%%%%%%%%%%%%%%%%%%%%%%%%%%%%%%%%%%%%%%%%%%%%%
\begin{frame}
	\frametitle{B2: LLC DC transformer}
	\begin{varblock}[\textwidth]{48\,V to 12\,V (4:1), 250\,W, about 1\,MHz}
		\begin{itemize}
			\item How is $L_\mathrm{m}$ sized for the charge $2Q_\mathrm{oss}+Q_\mathrm{TR}$ in the dead time ($Q_\mathrm{TR}$: charge of the transformer winding capacitance, L9)?
			\item Which winding arrangement, and what does it do to AC resistance and inter-winding capacitance (L7)?
			\item How is the synchronous rectifier (SR) timed when the package inductance falsifies the sensed drain voltage (L9)?
		\end{itemize}
	\end{varblock}
	\small
	Numbers to show: $L_\mathrm{m}$, dead time, gain at full and no load, transformer loss.
\end{frame}

%%%%%%%%%%%%%%%%%%%%%%%%%%%%%%%%%%%%%%%%%%%%%%%%%%%%%%%%%%%%%
%% Showcase B3 %%
%%%%%%%%%%%%%%%%%%%%%%%%%%%%%%%%%%%%%%%%%%%%%%%%%%%%%%%%%%%%%
\begin{frame}
	\frametitle{B3: Bidirectional DAB}
	\begin{varblock}[\textwidth]{48\,V to 12\,V and back, 200\,W, 0.5--1\,MHz}
		\begin{itemize}
			\item Is ZVS kept at light load and when $nV_2\neq V_1$ ($n$: turns ratio), and which modulation helps (L8)?
			\item Is the series inductance the transformer leakage or a separate part, and how accurate is it (L7)?
			\item Two switch nodes: how does the CM current flow through the transformer (L14)?
		\end{itemize}
	\end{varblock}
	\small
	Numbers to show: phase shift at full load, ZVS boundary, RMS current, series inductance.
\end{frame}

%%%%%%%%%%%%%%%%%%%%%%%%%%%%%%%%%%%%%%%%%%%%%%%%%%%%%%%%%%%%%
%% Showcase B4 %%
%%%%%%%%%%%%%%%%%%%%%%%%%%%%%%%%%%%%%%%%%%%%%%%%%%%%%%%%%%%%%
\begin{frame}
	\frametitle{B4: 6.78\,MHz wireless power}
	\begin{varblock}[\textwidth]{24\,V input, 15\,W into a specified coil and load range}
		\begin{itemize}
			\item What does the class E see when the coupling or the load changes, and how do you keep ZVS (L10)?
			\item Link efficiency at the worst coupling: which coil $Q$ do you need (L10)?
			\item Where do the harmonics of 6.78\,MHz fall, and what radiates (L14)?
		\end{itemize}
	\end{varblock}
	\small
	Numbers to show: $R$ and $C_\mathrm{sh}$ of the class E, drain peak voltage, coil $Q$, link efficiency.
\end{frame}

%%%%%%%%%%%%%%%%%%%%%%%%%%%%%%%%%%%%%%%%%%%%%%%%%%%%%%%%%%%%%
%% Showcase B5 %%
%%%%%%%%%%%%%%%%%%%%%%%%%%%%%%%%%%%%%%%%%%%%%%%%%%%%%%%%%%%%%
\begin{frame}
	\frametitle{B5: VHF dc--dc converter}
	\begin{varblock}[\textwidth]{12--24\,V to 5\,V, 5--10\,W, 13.56\,MHz}
		\begin{itemize}
			\item Which $C_\mathrm{oss}$ limits the frequency, and how much of it does the tank absorb (L10)?
			\item On/off control: modulation frequency, output ripple and what it does to the spectrum (L13, L14)?
			\item Air-core inductors: $Q$, stray field and the distance to the rest of the board (L7, L10)?
		\end{itemize}
	\end{varblock}
	\small
	Numbers to show: $C_\mathrm{oss}$ against $C_\mathrm{sh}$, drain peak voltage, on/off frequency, efficiency at 5 and 10\,W.
\end{frame}

%%%%%%%%%%%%%%%%%%%%%%%%%%%%%%%%%%%%%%%%%%%%%%%%%%%%%%%%%%%%%
%% Showcase B6 %%
%%%%%%%%%%%%%%%%%%%%%%%%%%%%%%%%%%%%%%%%%%%%%%%%%%%%%%%%%%%%%
\begin{frame}
	\frametitle{B6: Hybrid SC converter}
	\begin{varblock}[\textwidth]{48\,V to 12\,V (4:1), 150\,W, up to 1\,MHz}
		\begin{itemize}
			\item Which charge vectors, and which capacitance remains at DC bias (L12)?
			\item Is the charging soft, and how does the resonance move when the capacitance changes with voltage (L12)?
			\item How are the flying capacitors pre-charged at start-up?
		\end{itemize}
	\end{varblock}
	\small
	Numbers to show: $R_\mathrm{SSL}$ and $R_\mathrm{FSL}$, effective capacitance, resonant frequency, efficiency at full and half load.
\end{frame}

%%%%%%%%%%%%%%%%%%%%%%%%%%%%%%%%%%%%%%%%%%%%%%%%%%%%%%%%%%%%%
%% Evaluation %%
%%%%%%%%%%%%%%%%%%%%%%%%%%%%%%%%%%%%%%%%%%%%%%%%%%%%%%%%%%%%%
\subsection{Evaluation and outlook}

\begin{frame}
	\frametitle{Course evaluation}
	\begin{columns}[T]
		\begin{column}{0.55\textwidth}
			Please take ten minutes now for the course evaluation. This course runs for the first time; your answers decide what changes next year.
			\begin{itemize}
				\item Which lecture or tutorial helped you most in the design project?
				\item Where did the story line break for you?
				\item Notebooks, milestones, hardware track: what should stay, what should change?
			\end{itemize}
		\end{column}
		\begin{column}{0.4\textwidth}
			\begin{block}{Key message}
				Critical feedback is the most useful kind; please be specific.
			\end{block}
		\end{column}
	\end{columns}
\end{frame}

%%%%%%%%%%%%%%%%%%%%%%%%%%%%%%%%%%%%%%%%%%%%%%%%%%%%%%%%%%%%%
%% Theses %%
%%%%%%%%%%%%%%%%%%%%%%%%%%%%%%%%%%%%%%%%%%%%%%%%%%%%%%%%%%%%%
\begin{frame}
	\frametitle{Theses and projects at IAS}
	\begin{columns}[T]
		\begin{column}{0.48\textwidth}
			\small
			\begin{itemize}
				\item MHz GaN converters with planar and low-$\mu$ magnetics
				\item Wireless power at 6.78 and 13.56\,MHz
				\item Higher switching frequency for SiC converters in EV charging: the CLLC simulation model of our running example and the TI PMP41042 CLLLC test bench
				\item Data-driven core-loss models (MagNet) and their use in design
			\end{itemize}
		\end{column}
		\begin{column}{0.48\textwidth}
			\small
			\begin{itemize}
				\item Characterisation: $C_\mathrm{oss}$ loss, core loss and calorimetry at MHz
				\item Digital and ripple-based control of MHz converters
			\end{itemize}
			Promising designs from the hardware track can continue as a project, Studienarbeit or thesis. Talk to me after the session or by e-mail.
		\end{column}
	\end{columns}
\end{frame}

%%%%%%%%%%%%%%%%%%%%%%%%%%%%%%%%%%%%%%%%%%%%%%%%%%%%%%%%%%%%%
%% Key messages %%
%%%%%%%%%%%%%%%%%%%%%%%%%%%%%%%%%%%%%%%%%%%%%%%%%%%%%%%%%%%%%
\begin{frame}
	\frametitle{Key messages of the course}
	\begin{enumerate}
		\item $f_\mathrm{sw}$ is a design variable with an optimum: $L$ and $C$ shrink with $1/f_\mathrm{sw}$, losses, parasitics and delays do not.
		\item Hard switching costs $Q_\mathrm{oss}V_\mathrm{DC}$ per cycle; at MHz the energy must be recovered by soft switching, resonant or switched-capacitor circuits.
		\item At MHz the layout, the gate drive and the measurement are part of the circuit.
		\item Magnetics set the limit: core loss limits $\hat{B}$, skin and proximity effect limit the copper; low $\mu$, planar and air-core designs follow.
		\item Control, EMI and thermal design scale differently: delays stay, filters shrink, the heat density grows.
		\item A design is defended with numbers: loss budget, sensitivities and a validation plan.
	\end{enumerate}
\end{frame}

%%%%%%%%%%%%%%%%%%%%%%%%%%%%%%%%%%%%%%%%%%%%%%%%%%%%%%%%%%%%%
%% Exit ticket %%
%%%%%%%%%%%%%%%%%%%%%%%%%%%%%%%%%%%%%%%%%%%%%%%%%%%%%%%%%%%%%
% Two pages: questions (page 1) and answers (page 2).
\begin{frame}<handout:1-2>
	\frametitle{Exit questions}
	\begin{varblock}[\textwidth]{Three questions, one minute each}
		\begin{enumerate}
			\item Your design moves from 1 to 3\,MHz. Which three quantities do you recalculate first?
			\item Why does the EMI filter shrink faster with $f_\mathrm{sw}$ than the power inductor?
			\item Input and output power are measured with $\pm0.1$\,\% each at $\eta=98$\,\%. How uncertain is the loss?
		\end{enumerate}
	\end{varblock}
	\onslide<2-|handout:2>{
		\begin{enumerate}
			\item The switching energy per second ($Q_\mathrm{oss}V f_\mathrm{sw}$, or the ZVS charge in a shorter dead time), $\hat{B}$ and $F_\mathrm{R}$ of the magnetics, and the delays as a share of the period.
			\item The filter corner moves with $f_\mathrm{sw}$, so $L_\mathrm{f}$ and $C_\mathrm{f}$ fall with $1/f_\mathrm{sw}$. The power inductor carries the full ripple: core loss limits $\hat{B}$ and its VA per volume grows only with $f^{1-\alpha/\beta}$, e.g.\ $f^{0.44}$ (Tutorial 6).
			\item Loss $2\,\%$ of $P_\mathrm{in}$; the error $\pm(0.1+0.098)\,\%$ of $P_\mathrm{in}$ in the worst case is $\pm10\,\%$ of the loss. Measure the loss directly: calorimetry (L5).
		\end{enumerate}
	}
\end{frame}

%%%%%%%%%%%%%%%%%%%%%%%%%%%%%%%%%%%%%%%%%%%%%%%%%%%%%%%%%%%%%
%% References %%
%%%%%%%%%%%%%%%%%%%%%%%%%%%%%%%%%%%%%%%%%%%%%%%%%%%%%%%%%%%%%
\begin{frame}
	\frametitle{References}
	\scriptsize
	\begin{itemize}
		\item D.\,J. Perreault, ``Opportunities, progress and challenges in high-frequency power conversion,'' plenary talk, IEEE APEC 2024.
		\item A.\,J. Hanson, J.\,A. Belk, S. Lim, C.\,R. Sullivan, D.\,J. Perreault, ``Measurements and performance factor comparisons of magnetic materials at high frequency,'' IEEE Trans.\ Power Electron., vol.\ 31, no.\ 11, pp.\ 7909--7925, 2016.
		\item H. Zhang, \textit{Techniques for Efficient Wide-Range Radio-Frequency Power Generation}, PhD thesis, MIT, 2022.
		\item M. Joisher et al., ``High-performance high-power inductor design for high-frequency applications,'' IEEE APEC 2024.
		\item J.\,G. Kassakian, D.\,J. Perreault, G.\,C. Verghese, M.\,F. Schlecht, \textit{Principles of Power Electronics}, 2nd ed., Cambridge University Press, 2023.
		\item R.\,W. Erickson, D. Maksimovi\'c, \textit{Fundamentals of Power Electronics}, 3rd ed., Springer, 2020.
		\item Lectures L1 to L14 and Tutorials 1 to 13 of this course, with their references.
	\end{itemize}
\end{frame}
